% !TeX root = ../thuthesis-example.tex
\chapter{系统需求分析与方案论证}

\section{系统需求分析}

\subsection{功能需求分析}
结合项目当前实现，系统主要包含以下功能需求：
\begin{enumerate}
    \item 用户登录后能够进入统一主界面，并在多个功能页之间切换；
    \item 番茄计时模块支持任务名称输入、专注时长设置、休息时长设置、开始、暂停、继续和提前结束；
    \item 正向计时模块能够记录普通专注事件，并将统计结果写入本地数据文件；
    \item 日记模块能够按日期保存与读取文本内容，支持用户长期积累个人记录；
    \item 记账模块支持消费事项、金额和类型的录入，支持账单列表查看、编辑和删除；
    \item 统计模块能够对周、月等时间范围内的专注数据和消费数据进行汇总，展示图表和关键统计量；
    \item 系统应支持日间主题和夜间主题切换，以改善不同环境下的视觉体验。
\end{enumerate}
为了使系统具备更好的使用连续性，各模块之间还需要满足基本的数据联动要求。例如，番茄计时和正向计时产生的专注事件应自动进入本地记录文件，供日记模块展示和统计模块计算；记账模块录入的数据应当能够被统计模块实时读取并重新绘制图表；主题切换操作应对各页面和控件保持统一生效，避免界面风格割裂。由此可见，本系统的功能需求并非简单的模块堆叠，而是强调功能之间的协同工作能力。

\subsection{非功能需求分析}
在非功能层面，系统需要满足以下要求：首先，界面交互应清晰直观，便于普通用户快速上手；其次，系统运行应保持较好的响应性，避免在计时刷新、文件读写和图表绘制过程中出现明显卡顿；再次，数据存储应以本地化为主，降低部署与使用门槛；最后，系统结构应具备较好的模块化特征，便于后续扩展同步功能、数据库功能和更多工具模块。
除此之外，系统还需要满足一定的可维护性和可扩展性要求。由于毕业设计项目往往会在后续继续迭代，因此代码结构应尽量清晰，页面类、服务类和数据结构之间的职责划分应当明确。以当前项目为例，用户会话由独立的Session对象维护，数据读写由DataStore、DiaryStore和AccountingStore等服务类封装，统计逻辑则由StatsService和AccountingStatsService单独负责，这种分工方式有利于减少业务耦合度，并为后续引入数据库、云同步或更复杂的分析模型提供基础。

\section{关键技术与方案比较}

\subsection{GUI框架选型}
在桌面图形界面技术的选择上，本文对Qt、Electron和WPF进行了比较。Electron开发效率较高，但资源占用相对较大；WPF在Windows平台具有较好的开发体验，但跨平台能力有限；Qt则兼顾了性能、跨平台能力和组件丰富度，特别适合毕业设计类桌面系统开发。因此，本文最终选择Qt Widgets作为界面开发框架。
具体而言，Qt Widgets 提供了QMainWindow、QTabWidget、QTableWidget、QMenu、QToolBar、QCalendarWidget等大量成熟控件，能够快速搭建功能型桌面应用界面。更重要的是，Qt拥有稳定的信号与槽机制，便于描述页面内外的交互关系。例如，登录页面可通过发出登录成功信号进入主窗口，记账页面可通过动作信号响应编辑和删除操作，统计页面可通过按钮与图表点击事件完成图形切换与详情展示。这些特性使Qt在本项目中具备较高的开发效率与实现灵活性。

\subsection{数据存储方案比较}
在数据存储方案方面，可选方案主要包括关系数据库、本地文本文件和JSON文件。考虑到本系统定位为轻量级个人桌面工具，数据规模有限、部署环境简单，若直接引入数据库会增加配置与维护成本。因此，系统最终采用混合本地文件存储策略：专注记录与日记内容使用文本文件保存，便于按日期管理；记账数据则采用JSON对象逐行存储，既保留结构化特征，又降低解析复杂度。
该方案还有一个显著优点，即可视化和调试过程更加直接。开发者可以直接打开对应用户目录下的文件查看原始数据，有利于快速定位问题和验证功能效果。对于毕业设计项目而言，这种方式比完整引入数据库系统更适合当前规模，也更符合“轻量、直观、可演示”的实现目标。当然，这种方案的局限性也较为明显，例如复杂查询能力有限、缺少事务和并发控制，因此在后续展望部分，本文也提出了引入SQLite等数据库的改进方向。

\subsection{图表实现方案比较}
图表功能可通过引入第三方图表库实现，也可基于Qt绘图机制自行开发。第三方库开发速度较快，但会增加项目依赖；自定义绘制虽然实现复杂度较高，但更容易控制样式、交互和主题适配。结合本项目规模和教学实践要求，系统最终采用自定义饼状图和柱状图组件实现数据可视化，以增强对绘制流程和事件交互的掌握。
从系统效果来看，自定义图表组件不仅能够绘制静态图形，还能够与页面交互逻辑结合。例如，统计页面可以根据时间范围切换不同数据源，并将汇总结果同步展示到饼状图和柱状图中；用户点击图形元素后，还可以触发对应的数据详情展示。这种“绘制+交互”的设计方式，更有利于体现项目的工程实现能力，也使系统在展示环节更具完整性。

\section{系统总体方案确定}

\subsection{架构设计思路}
系统采用模块化设计思想，以主窗口作为统一容器，通过选项卡组织各业务页面。业务逻辑上，界面层负责输入输出与交互反馈，服务层负责数据读写和统计计算，会话层负责维护当前登录用户状态。该设计能够较好地降低模块耦合度，提高后续维护和扩展效率。
从程序启动流程来看，系统首先创建QApplication对象并加载已保存的主题配置，然后进入登录窗口。登录成功后，系统创建主窗口，并将番茄钟、正向计时、数据统计、记账、日记等核心功能页面注册到标签页中。通过这种设计，项目实现了较为清晰的“入口窗口—主业务窗口—具体功能页面”层级结构，有利于控制页面职责边界并统一管理全局外观和用户状态。

\subsection{模块划分方案}
根据系统功能，本文将项目划分为登录会话模块、番茄计时模块、正向计时模块、日记模块、记账模块、数据统计模块、主题管理模块和扩展工具模块。其中，记账模块与统计模块是本次系统完善中的核心新增内容，它们直接提升了系统的综合应用价值与展示效果。
各模块之间通过较为明确的接口进行协作。例如，会话模块为其他页面提供当前用户名；数据存储模块为业务页面提供持久化能力；统计模块在需要时主动读取存储层数据并形成汇总结果；主题管理模块则通过统一样式表影响整个应用的视觉表现。这种划分方式不仅提升了系统的可理解性，也有助于在论文中清晰描述各部分职责。